% Part: computability
% Chapter: recursive-functions
% Section: pr-functions

\documentclass[../../../include/open-logic-section]{subfiles}

\begin{document}

\olfileid[pt]{cmp}{rec}{prf}
\olsection{Funções recursivas primitivas}

Registremos novamente como podemos definir novas funções a partir de funções
existentes usando recursão primitiva e composição.

\begin{defn}
  \ollabel{defn:primitive-recursion} Suponha que $f$ seja uma função de
  $k$ argumentos ($k\ge 1$) e que $g$ seja uma função de $(k+2)$ argumentos.
  A função definida por \emph{recursão primitiva a partir de $f$ e $g$} é a
  função de $(k+1)$ argumentos~$h$ definida pelas equações
  \begin{align*}
    h(x_0,\dots,x_{k-1},0) & =  f(x_0,\dots,x_{k-1}) \\
    h(x_0,\dots,x_{k-1},y+1) & = g(x_0,\dots,x_{k-1}, y, h(x_0,\dots,x_{k-1}, y))
  \end{align*}
\end{defn}

\begin{defn}
  \ollabel{defn:composition} Suponha que $f$ seja uma função de $k$ argumentos
  e que $g_0$, \dots, $g_{k-1}$ sejam $k$ funções, todas de $n$ argumentos.
  A função definida por \emph{composição a partir de $f$ e de $g_0$,
    \dots,~$g_{k-1}$} é a função de $n$ argumentos~$h$ definida por
  \[
  h(x_0, \dots, x_{n-1}) =
  f(g_0(x_0, \dots, x_{n-1}), \dots, g_{k-1}(x_0, \dots, x_{n-1})).
  \]
\end{defn}

Além de $\Succ$ e das funções de projeção
\[
\Proj{n}{i}(x_0,\dots,x_{n-1}) = x_i,
\]
para cada número natural $n$ e cada $i < n$, incluiremos entre as funções
recursivas primitivas a função~$\Zero(x) = 0$.

\begin{defn}
  O conjunto das funções recursivas primitivas é o conjunto de funções de
  $\Nat^n$ em $\Nat$ definido indutivamente pelas cláusulas seguintes:
  \begin{enumerate}
  \item $\Zero$ é recursiva primitiva.
  \item $\Succ$ é recursiva primitiva.
  \item Cada função de projeção $\Proj{n}{i}$ é recursiva primitiva.
  \item Se $f$ é uma função recursiva primitiva de $k$ argumentos e $g_0$,
    \dots,~$g_{k-1}$ são funções recursivas primitivas de $n$ argumentos,
    então a composição de $f$ com $g_0$, \dots,~$g_{k-1}$ é recursiva
    primitiva.
  \item Se $f$ é uma função recursiva primitiva de $k$ argumentos e $g$ é
    uma função recursiva primitiva de $k+2$ argumentos, então a função
    definida por recursão primitiva a partir de $f$ e $g$ é recursiva
    primitiva.
\end{enumerate}
\end{defn}

\begin{explain}
De modo mais conciso, o conjunto das funções recursivas primitivas é o menor
conjunto que contém $\Zero$, $\Succ$ e as funções de
projeção~$\Proj{n}{j}$ e que é fechado sob composição e recursão primitiva.

Outra maneira de descrever o conjunto das funções recursivas primitivas é
defini-lo em termos de ``estágios''. Seja $S_0$ o conjunto das funções
iniciais: $\Zero$, $\Succ$ e as projeções. Essas são as funções recursivas
primitivas do estágio~$0$. Uma vez definido um estágio $S_i$, seja $S_{i+1}$
o conjunto de todas as funções obtidas pela aplicação de uma única instância
de composição ou de recursão primitiva a funções que já pertencem a~$S_i$.
Então
\[
S = \bigcup_{i \in \Nat} S_i
\]
é o conjunto de todas as funções recursivas primitivas.
\end{explain}

Verifiquemos que $\Add$ é uma função recursiva primitiva.

\begin{prop}
  A função adição $\Add(x,y) = x+y$ é recursiva primitiva.
\end{prop}

\begin{proof}
Já temos uma definição por recursão primitiva de $\Add$ em termos de duas
funções $f$ e~$g$ que corresponde ao formato de
\olref{defn:primitive-recursion}:
\begin{align*}
  \Add(x_0, 0) & = f(x_0) = x_0\\
  \Add(x_0, y+1) & = g(x_0, y, \Add(x_0, y)) =
  \Succ(\Add(x_0, y))
\end{align*}
Assim, $\Add$ é recursiva primitiva desde que $f$ e $g$ também o sejam.
$f(x_0) = x_0 = \Proj{1}{0}(x_0)$, e as funções de projeção contam como
recursivas primitivas; logo, $f$ é recursiva primitiva. A função $g$ é a
função de três argumentos $g(x_0, y, z)$ definida por
\[
g(x_0, y, z) = \Succ(z).
\]
Isso ainda não mostra que $g$ seja recursiva primitiva, pois $g$ e $\Succ$
não são exatamente a mesma função: $\Succ$ é unária, e $g$ precisa ter três
argumentos. Mas podemos definir $g$ ``oficialmente'' por composição como
\[
g(x_0, y, z) = \Succ(\Proj{3}{2}(x_0, y, z))
\]
Como $\Succ$ e $\Proj{3}{2}$ contam como funções recursivas primitivas, $g$
também conta, pois pode ser definida por composição a partir de funções
recursivas primitivas.
\end{proof}

\begin{prop}
  \ollabel{prop:mult-pr}
  A função multiplicação $\Mult(x,y) = x \cdot y$ é recursiva primitiva.
\end{prop}

\begin{proof}
  Exercício.
\end{proof}

\begin{prob}
  Demonstre \olref[cmp][rec][prf]{prop:mult-pr} mostrando que a definição por
  recursão primitiva de $\Mult$ pode ser posta na forma exigida por
  \olref[cmp][rec][prf]{defn:primitive-recursion} e que as funções
  correspondentes $f$ e~$g$ são recursivas primitivas.
\end{prob}

\begin{ex}
Eis nosso primeiro exemplo de uma definição por recursão primitiva:
\begin{align*}
h(0) & =  1 \\
h(y+1) & =  2 \cdot h(y).
\end{align*}
Essa função não pode se ajustar à forma exigida por
\olref{defn:primitive-recursion}, pois $k=0$. A definição também envolve as
constantes $1$ e $2$. Para contornar o primeiro problema, introduzamos um
argumento fictício e definamos a função~$h'$:
\begin{align*}
h'(x_0, 0) & =  f(x_0) = 1 \\
h'(x_0, y+1) & =  g(x_0, y, h'(x_0, y)) = 2 \cdot h'(x_0, y).
\end{align*}
A função $f(x_0) = 1$ pode ser definida a partir de $\Succ$ e $\Zero$ por
composição: $f(x_0) = \Succ(\Zero(x_0))$. A função $g$ pode ser definida por
composição a partir de $g'(z) = 2 \cdot z$ e de projeções:
\begin{align*}
g(x_0, y, z) & = g'(\Proj{3}{2}(x_0, y, z))
\intertext{e $g'$, por sua vez, pode ser definida por composição como}
g'(z) & = \Mult(g''(z), \Proj{1}{0}(z))
\intertext{e}
g''(z) & = \Succ(f(z)),
\end{align*}
onde $f$ é como acima: $f(z) = \Succ(\Zero(z))$. Agora que temos~$h'$, podemos
usar novamente a composição e tomar
$h(y) = h'(\Proj{1}{0}(y),\Proj{1}{0}(y))$. Isso mostra que $h$ pode ser
definida a partir das funções básicas usando uma sequência de composições e
recursões primitivas; portanto, $h$ é recursiva primitiva.
\end{ex}

\end{document}
